\chapter{Los números decimales}\label{ch:g5:decimals}

Las décimas y las centésimas aparecieron en el \cref{ch:g4:decimals}; añadiendo un
nivel más, las milésimas, se completa el sistema decimal. Este capítulo
consolida la lectura, la descomposición y --- sobre todo --- la comparación de
números decimales sin caer en sus famosas trampas.

\section{Hasta las milésimas}

\begin{definition}[Las cifras decimales]\label{def:g5:decimals:places}
Después del \omterm{def:g4:decimals:point}{punto decimal} vienen las décimas, las centésimas y las
milésimas --- cada una vale diez veces menos que la anterior:
\[
4.362 = 4 + \frac{3}{10} + \frac{6}{100} + \frac{2}{1000}
= \frac{4\,362}{1000}.
\]
\end{definition}

\begin{omfigure}
\begin{tikzpicture}[scale=0.9]
  \draw (0,0) grid[xstep=2.3, ystep=0.85] (9.2,1.7);
  \node[font=\small] at (1.15,1.275) {unidades};
  \node[font=\small] at (3.45,1.275) {décimas};
  \node[font=\small] at (5.75,1.275) {centésimas};
  \node[font=\small] at (8.05,1.275) {milésimas};
  \node[font=\large, omDef] at (1.15,0.425) {$4$};
  \node[font=\large, omThm] at (3.45,0.425) {$3$};
  \node[font=\large, omThm] at (5.75,0.425) {$6$};
  \node[font=\large, omThm] at (8.05,0.425) {$2$};
  \fill (2.3,0.06) circle (2.5pt);
\end{tikzpicture}

{\small La tabla de valor posicional del $4.362$: parte entera en azul y parte
decimal en rojo, con el punto entre las unidades y las décimas.}
\end{omfigure}

\begin{example}[Muchas lecturas, un solo número]\label{ex:g5:decimals:readings}
$4.362$ se puede leer «cuatro punto tres seis dos», o «cuatro y
$362$ milésimas», o descomponerse como $4 + 0.3 + 0.06 + 0.002$. Y
añadir \omterm{def:g1:counting:zero}{ceros} al final no cambia nada: $4.362 = 4.3620$. Pero
$4.362 \neq 4.0362$: un \omterm{def:g1:counting:zero}{cero} \emph{intercalado} empuja cada cifra a un
lugar menor.
\end{example}

\begin{example}[¿Dónde está en la recta?]\label{ex:g5:decimals:line}
Para situar $4.362$: entre $4$ y $5$; ampliando, entre $4.3$ y
$4.4$; ampliando otra vez, entre $4.36$ y $4.37$, más cerca de $4.36$.
Cada cifra decimal es un nivel más de ampliación.
\end{example}

\begin{omfigure}
\begin{tikzpicture}[scale=1.15]
  \draw[-Stealth] (-0.2,1.6) -- (10.4,1.6);
  \foreach \i in {0,...,10} \draw (\i,1.52) -- (\i,1.68);
  \node[below, font=\small] at (0,1.5) {$4$};
  \node[below, font=\small] at (10,1.5) {$5$};
  \node[below, font=\small] at (3,1.5) {$4.3$};
  \node[below, font=\small] at (4,1.5) {$4.4$};
  \fill[omThm] (3.62,1.6) circle (1.7pt);
  \draw[omExo, thin] (3,1.45) -- (0,0.15);
  \draw[omExo, thin] (4,1.45) -- (10,0.15);
  \draw[-Stealth] (-0.2,0) -- (10.4,0);
  \foreach \i in {0,...,10} \draw (\i,-0.08) -- (\i,0.08);
  \node[below, font=\small] at (0,-0.1) {$4.3$};
  \node[below, font=\small] at (10,-0.1) {$4.4$};
  \node[below, font=\small] at (6,-0.1) {$4.36$};
  \fill[omThm] (6.2,0) circle (1.7pt)
    node[above, font=\small] {$4.362$};
\end{tikzpicture}

{\small Ampliación de la recta numérica: el tramo de $4.3$ a $4.4$,
aumentado, está a su vez cortado en diez centésimas --- y $4.362$ queda
justo pasado el $4.36$.}
\end{omfigure}

\section{Comparar decimales}

\begin{method}[Comparar dos decimales]\label{met:g5:decimals:compare}
\begin{enumerate}
  \item compara las partes enteras: $6.1 > 5.987$;
  \item con la misma parte entera: compara las partes decimales cifra a cifra
        desde el punto --- décimas, luego centésimas y luego
        milésimas;
  \item completar con \omterm{def:g1:counting:zero}{ceros} finales hace justa la comparación:
        $2.5$ frente a $2.48$ pasa a ser $2.50$ frente a $2.48$, y $50 > 48$
        centésimas.
\end{enumerate}
La trampa, una vez más: \emph{una parte decimal más larga no significa un
número mayor} --- $2.48$ tiene más cifras que $2.5$ y es menor.
\end{method}

\begin{example}\label{ex:g5:decimals:order}
Ordena $7.3$; $7.09$; $7.31$; $7.299$. Completa a tres decimales:
$7.300$; $7.090$; $7.310$; $7.299$. Entonces,
\[
7.09 < 7.299 < 7.3 < 7.31 .
\]
Fíjate en que $7.299 < 7.3$ aunque $299$ parezca grande: son $299$ milésimas
frente a $300$ milésimas.
\end{example}

\begin{example}[Colar un número entre dos decimales]\label{ex:g5:decimals:between}
¿Hay algún número entre $5.7$ y $5.8$? Sí, muchísimos: $5.75$,
$5.71$, $5.799$\,\dots\ ¿Y entre $5.79$ y $5.8$? También muchos:
$5.795$, por ejemplo. Entre dos decimales distintos \emph{siempre} se puede
colar otro --- no existe el decimal «siguiente» a
uno dado.
\end{example}

\section{Redondear decimales}

\begin{definition}[Redondeo]\label{def:g5:decimals:round}
Para redondear a la unidad (o a la décima, o a la centésima), mira la cifra
siguiente: con $5$ o más se redondea hacia arriba y, si no, hacia abajo. Así, $4.362$
redondea a $4$ (unidad), a $4.4$ (décima) y a $4.36$ (centésima).
\end{definition}

\begin{example}[El dinero se redondea a los céntimos]\label{ex:g5:decimals:money}
Un precio calculado como $7.4983$ se muestra como $7.50$: las cantidades
de la vida real se redondean a la centésima. Fíjate en la cascada: $7.4983 \to
7.50$, donde el $49$ pasó a $50$ --- el \omterm{def:g4:numbers:round}{redondeo} puede propagarse por
varias cifras.
\end{example}

\section{Ejercicios}

\begin{exercise}[$\star$]\label{exo:g5:decimals:1}
Escribe como número decimal: $5 + \frac{2}{10} + \frac{7}{1000}$;
\quad $\frac{85}{100}$; \quad $\frac{4\,507}{1000}$; \quad «doce
y nueve milésimas».
\end{exercise}

\begin{exercise}[$\star$]\label{exo:g5:decimals:2}
En $27.418$: ¿qué cuenta el $4$? ¿Y el $8$? Descompón el número
como en la \cref{def:g5:decimals:places}.
\end{exercise}

\begin{exercise}[$\star$]\label{exo:g5:decimals:3}
¿Verdadero o falso? Explícalo con los lugares.
$3.60 = 3.6$; \quad $0.5 = 0.05$; \quad $2.070 = 2.07$; \quad
$1.3 = 1.03$.
\end{exercise}

\begin{exercise}[$\star$]\label{exo:g5:decimals:4}
Copia y completa con $<$, $>$ o $=$:
\[
4.8 \;?\; 4.79, \qquad
0.302 \;?\; 0.32, \qquad
6.5 \;?\; 6.500, \qquad
9.09 \;?\; 9.1 .
\]
\end{exercise}

\begin{exercise}[$\star$]\label{exo:g5:decimals:5}
Ordena de menor a mayor: $3.14$; \ $3.4$; \ $3.041$; \
$3.104$; \ $3.41$.
\end{exercise}

\begin{exercise}[$\star$]\label{exo:g5:decimals:6}
¿Entre qué dos números enteros está cada número? ¿Y entre qué
dos números de una cifra decimal? \quad $6.83$; \quad $0.492$; \quad
$12.06$.
\end{exercise}

\begin{exercise}[$\star$]\label{exo:g5:decimals:7}
Redondea $8.276$: a la unidad; a la décima; a la centésima. Redondea
$3.95$ a la décima.
\end{exercise}

\begin{exercise}[$\star$]\label{exo:g5:decimals:8}
Da un número estrictamente comprendido entre: $2.6$ y $2.7$; \quad $0.41$ y
$0.42$; \quad $5.99$ y $6$.
\end{exercise}

\begin{exercise}[$\star$]\label{exo:g5:decimals:9}
Las alturas de cuatro plantas son $0.85$ m, $1.2$ m, $0.9$ m y
$1.02$ m. Ordénalas de la más baja a la más alta.
\end{exercise}

\begin{exercise}[$\star\star$]\label{exo:g5:decimals:10}
Lucía dice: «$7.12 > 7.9$ porque $12 > 9$». Corrige su error
comparando primero las décimas.
\end{exercise}

\begin{exercise}[$\star\star$]\label{exo:g5:decimals:11}
Halla todos los números con exactamente dos cifras decimales que redondean a
$6.4$ al redondear a la décima. ¿Cuáles son el menor y el
mayor?
\end{exercise}
